\section{Introduction}
\label{sec:introduction}

% 1. scirm 的 introduction 写法
% - 先介绍了科学写作的生成很多，但是不好评估
% - 然后介绍了为每个科学写作微调成本高，少资源
% - 再引入 LLm-as-a-judge 是个方法，但是缺少领域知识，已有的在训练时引入
%   但是又被固化了，推理时还是不够遵循不同维度、准则等等
% - 最后引入他们的方法，奖励推理模型。先说了已有的奖励推理模型的几个不足：
%   首先是只在通用评价任务上，没有科学写作评价；另一个是采用成对偏好建模，对
%   文本质量打分不够；还有就是训练时候适配了固定打分规则，换个标准不同的任务就不
%   行；最后就是也有论文做了专业领域打分，但仅适用单一任务，且依赖闭源商业大模型，
%   拓展泛化能力不足。
% - 最后介绍了他们的方法：两阶段推理训练。以及介绍了好处：就是可以学习既定的评定
%   准则，同时还能反向解读准则，依次修正自身推理，强调这真是现在的奖励推理模型没有的。

Large language models (LLMs) are increasingly used to score scientific text
against task-specific criteria. A common design asks the judge to generate a
rationale before returning its score, with the expectation that explicit
reasoning will support a more reliable judgment. Yet the available evidence
does not show a uniform scoring benefit. Explanation requests can change
agreement with human ratings, and chain-of-thought (CoT) can even degrade the
judgment distribution of an LLM judge
\citep{chiang-lee-2023-closer,wang-etal-2025-improving-llm-judge}. Likewise,
rationale-augmented fine-tuning has improved some tasks but trailed label-only
training on others
\citep{hsieh-etal-2023-distilling,shi-etal-2025-rationales-nlu,
su-etal-2026-rationale-disease}. Thus, producing a plausible rationale does not
by itself establish that the associated score is learned more reliably.

This inconsistency motivates a more precise question: when a judge outputs
both a rationale and a score, which part of the pipeline accounts for any
change in scoring? Rationale-first inference asks a fixed model to reason
before answering, whereas rationale supervision also optimizes the rationale
tokens during fine-tuning. Conflating the two can obscure whether a score
changes because the model learned a better criterion-specific decision, or
because it became adapted to an inference interface that expects reasoning.
We therefore treat rationale supervision and rationale-first inference as
separate experimental factors rather than as a single CoT intervention.

Existing generative judges do not provide this separation. Prometheus learns
rubric-conditioned feedback followed by a score, and SciRM trains reasoning
reward models for multiple scientific-writing criteria
\citep{kim-etal-2024-prometheus-iclr,sahinuc-etal-2026-reward}. More closely,
work on natural language understanding compares label-only and
rationale-augmented fine-tuning and proposes Align to balance paired direct and
reasoning views \citep{shi-etal-2025-rationales-nlu}. However, its evidence
comes from general understanding tasks rather than rubric-conditioned
scientific-writing evaluation. Whether the same supervision choices improve
scientific-writing scores, and how they interact with the inference format,
therefore remains unresolved.

In this work, we conduct a controlled study of whether CoT supervision improves
scientific-writing judgments under explicit rubrics. Concretely, we train
matched
judges with label-only targets, CoT targets, or Align, which pairs direct and
reasoning supervision. Rather than treating CoT as a single design choice, we
independently vary the training target and the inference interface. We apply
this design to seven tasks spanning related-work evaluation and review utility,
each with an explicit criterion and label space. This separation allows us to
identify whether a scoring change reflects supervision, reliance on a
reasoning-first interface, or their interaction. We further trace paired
prediction changes to distinguish invalid output schemas from shifts among
valid scores.

Our contributions and findings are threefold. First, we systematically compare
label-only and CoT supervision across multiple scientific-writing criteria and
show that CoT is not a uniform improvement: its effects vary by task and
inference interface. Second, we demonstrate that performance degradation has
two distinct forms, output-schema failure and valid-score drift, which are
obscured by aggregate metrics. Third, we adapt Align to preserve direct
supervision while learning a reasoning path. Align recovers roughly three
quarters of the label-only-correct predictions lost after CoT fine-tuning
across two seeds, although our current evidence does not attribute this
recovery solely to its loss-balancing design. Together, these findings show
that rationale generation and scoring reliability must be evaluated as
related but distinct objectives.
